\documentclass[12pt]{article}
\usepackage[utf8]{inputenc}
\usepackage[T1]{fontenc}
\usepackage{hyperref}
\title{Mur Pano Language Resources for K-12 Curriculum Development}
\author{Compiled by Claude Sonnet 5 (claude-sonnet-5) for Andr\'as Kornai}
\date{}
\begin{document}
\maketitle

\section{Introduction}
Mur Pano (tkv) is an Austronesian language of Papua New Guinea. Wikipedia places it in Mur village on the north coast of Madang Province and gives 700--800 native speakers. Joshua Project places the Mur Pano people in East Sepik Province and gives a population of 1,500. The sources conflict on location. The language is distinct from Murupi (mqw), whose resources are not used here.

\section{Primary Educational Resources}
None were found.

\section{Linguistic and Reference Materials}
The only source found is the Wikipedia article. \url{https://en.wikipedia.org/wiki/Mur_Pano_language}

\section{Religious and Community Resources}
Joshua Project has a people-group page for the Mur Pano. \url{https://joshuaproject.net/people_groups/20999/PP}

\section{Digital and Audio Resources}
Global Recordings Network has a language page for tkv. I did not verify what recordings it holds. \url{https://globalrecordings.net/en/language/tkv}

\section{Curriculum Development Context}
The language is not known to be taught in schools. Documented material is very thin. A curriculum would need new documentation with the community.
\end{document}